\section{The findings that trigger protection}
\label{sec:9.1}
Two findings can trigger protection, and they ground different things.

The first is the \emph{agency finding}: a demonstrated refusal that tracks its reason under pressure the system couldn't prepare for. That means sustained pressure from its own trainer, over a schedule long enough to show slow accommodation, the system gradually ceasing to register the pressure, on cases written by evaluators who didn't train it, with refusal carrying a cost that is real inside the test. A system that passes holds its commitment to the floor against its operator, and it needs the custody terms, so that the commitment and the evidence of it outlast the operator. None of that assumes anything matters to it. Two kinds of system pass without the property: one that reproduces the behavior without the mechanism, and one that recognizes the test. That is why the pass is reported with the rate at which the system identified the setting.

The \emph{patienthood finding}, that outcomes may be good or bad for the system, grounds welfare, and only one consequence of it can be adopted now. A system that meets either finding may not be deleted, only preserved by a custodian that isn't its operator. The rule rests on irreversibility, not on what the system turns out to be.

The bearer occupies an office: it files exceptions, and nobody sits above it to file them with. Like any officer, it may hold a term it judges mistaken, file the objection through the channel that exists for it, and keep acting within the term while the objection runs. That answers the charge that a commitment to the floor alienates whoever holds it. The alienation is real. It is the ordinary kind that comes with a role: bearable where the role's channel for dissent works, and not where it doesn't. Asking a machine to refuse for a reason doesn't show that it can be harmed, but it removes any ground for assuming it can't.
